\section{Pregled literature}

\subsection{Pregled Arhitekture}
Implementacija chatbot sustava za studentska pitanja temelji se na modularnoj arhitekturi koja omogućava skalabilnost i održivost. Sustav je dizajniran kao mikroservisna aplikacija koja integrira semantic routing kao osnovni mehanizam za usmjeravanje korisničkih upita prema relevantnim odgovorima.
Arhitektura sustava sastoji se od pet glavnih komponenti: korisničko sučelje (frontend), API gateway, semantic routing engine, baza znanja i recommendation servis. Svaka komponenta ima jasno definirane odgovornosti i komunicira s ostalima kroz standardizirane API-je.

\subsection{Tehnološki Stack}
Odabir tehnologija temeljio se na potrebi za visokom performansom, skalabilnošću i jednostavnošću održavanja. Backend je implementiran u Python programskom jeziku koristeći FastAPI framework zbog njegove brzine i automatske dokumentacije API-ja.
Za embedding modele korišten je sentence-transformers s pretreniranim modelom 'all-MiniLM-L6-v2' koji pruža dobar balans između kvalitete vektora i brzine izvršavanja. Vector store implementiran je pomoću FAISS biblioteke koja omogućava brzo pretraživanje najbližih susjeda u višedimenzionalnom prostoru.

\subsection{Motivacija}

Motivacija za razvoj Sembota proizlazi iz potrebe za automatizacijom podrške studentima, posebno u kontekstu sveučilišta s velikim brojem upita. Projekt ima za cilj smanjiti opterećenje na administrativno osoblje, omogućiti studentima brži pristup informacijama i stvoriti platformu koja se može lako prilagoditi različitim fakultetima i obrazovnim institucijama.

\subsection{Ciljevi projekta}

\begin{enumerate}
    \item Razviti chatbot koji može odgovarati na česta pitanja studenata koristeći semantičko pretraživanje.
    \item Omogućiti jednostavno proširenje sustava dodavanjem novih pitanja i odgovora.
    \item Osigurati skalabilnost i fleksibilnost sustava za buduće nadogradnje.
    \item Pružiti intuitivno korisničko sučelje dostupno putem web preglednika.
\end{enumerate}

\subsection{Problem}
Studenti se suočavaju s brojnim izazovima prilikom traženja informacija o fakultetskim procesima. Na primjer, pitanja poput "Koji su uvjeti za upis?" ili "Kada počinju predavanja?" često zahtijevaju trenutne odgovore, ali administrativni uredi nemaju dovoljno resursa za obradu velikog broja upita. Osim toga, postojeći sustavi, poput FAQ stranica, često su ograničeni na unaprijed definirane odgovore i ne mogu obraditi varijacije u formulaciji pitanja. Ovo dovodi do neefikasnosti, povećanog opterećenja na osoblje i lošijeg korisničkog iskustva.

Ključni izazovi:
\begin{enumerate}
    \item Veliki broj upita: Sveučilišta dobivaju stotine ili tisuće upita dnevno tijekom upisnih rokova.
    \item Raznolikost formulacija: Studenti postavljaju ista pitanja na različite načine, što otežava automatsku obradu.
    \item Ograničeni resursi: Administrativno osoblje ne može uvijek pružiti trenutne odgovore.
    \item Skalabilnost: Postojeći sustavi često nisu prilagođeni za obradu velikih količina podataka ili novih vrsta upita.
\end{enumerate}

\subsection{Alternativna rješenja}

Postoji nekoliko pristupa za rješavanje ovog problema, svaki s prednostima i nedostacima:

\begin{enumerate}
   \item \textbf{Statične FAQ stranice}:
   \begin{itemize}
       \item \textbf{Opis}: Web stranice s unaprijed definiranim pitanjima i odgovorima.
       \item \textbf{Prednosti}: Jednostavne za implementaciju, ne zahtijevaju složene tehnologije.
       \item \textbf{Nedostaci}: Ne mogu obraditi varijacije u formulaciji pitanja, nisu interaktivne, teško se skaliraju za veliki broj pitanja.
       \item \textbf{Primjena}: Pogodne za male institucije s ograničenim brojem upita.
   \end{itemize}
   
   \item \textbf{E-mail ili telefonska podrška}:
   \begin{itemize}
       \item \textbf{Opis}: Studenti kontaktiraju administrativno osoblje putem e-pošte ili telefona.
       \item \textbf{Prednosti}: Omogućuje personalizirane odgovore, pogodno za složene upite.
       \item \textbf{Nedostaci}: Vremenski zahtjevno, ne skalira se dobro s velikim brojem upita, skupo za održavanje.
       \item \textbf{Primjena}: Prikladno za specifične, rijetke upite koji zahtijevaju ljudsku intervenciju.
   \end{itemize}
   
   \item \textbf{Chatbotovi temeljeni na ključnim riječima}:
   \begin{itemize}
       \item \textbf{Opis}: Sustavi koji traže podudarnosti ključnih riječi između upita i baze podataka.
       \item \textbf{Prednosti}: Relativno jednostavni za implementaciju, brzi za male baze podataka.
       \item \textbf{Nedostaci}: Ne razumiju kontekst ili semantičku sličnost, što dovodi do netočnih odgovora za varijacije pitanja.
       \item \textbf{Primjena}: Pogodni za jednostavne sustave s ograničenim brojem pitanja.
   \end{itemize}
   
   \item \textbf{AI chatbotovi s prirodnim jezikom}:
   \begin{itemize}
       \item \textbf{Opis}: Chatbotovi temeljeni na naprednim modelima poput BERT-a, GPT-a ili sličnih.
       \item \textbf{Prednosti}: Razumiju prirodni jezik, mogu obraditi složene upite, prilagodljivi su.
       \item \textbf{Nedostaci}: Zahtijevaju značajne računalne resurse, složeni su za implementaciju, često zahtijevaju plaćene API-je.
       \item \textbf{Primjena}: Idealni za velike institucije s velikim brojem upita i dostupnim resursima.
   \end{itemize}
\end{enumerate}

Sembot kombinira prednosti AI chatbotova s jednostavnošću i skalabilnošću, koristeći Semantic Router za razumijevanje prirodnog jezika i FastAPI za efikasnu obradu upita. Ovaj pristup omogućuje fleksibilnost i lako proširenje sustava uz minimalne resurse.


